\section{空间的性质}
\label{sec:topo-properties}

前几节的例子显示，拓扑空间可以离日常的几何直觉很远：平凡拓扑中所有点都挤在一起，点列同时收敛到每个点；余可数拓扑中点列几乎探测不到闭包。本节集中研究整个空间满足的几类条件，它们排除这些病态现象，或者保证局部信息能够汇成整体信息。分离性问不同点能否被邻域分开；可数性与可分性问开集、邻域和稠密点集能否用可数个代表描述；连通性问空间能否拆成两块；紧性问开覆盖能否化为有限个。每项条件都用前面的点列、映射与构造来检验它的后果。

\subsection{分离性}
\label{sec:topo-separation}

在$\mathbb R$中，两个不同的点$x<y$可以分别放进$(x-r,x+r)$与$(y-r,y+r)$，其中$r=(y-x)/2$，两个区间不相交。平凡拓扑中做不到这一点，因为每个点的唯一开邻域都是全空间；正是这个缺陷让点列在平凡拓扑中有多个极限。

\begin{definition}[Hausdorff空间]
\label{def:topo-hausdorff}
若对$X$中任意两个不同的点$x\neq y$，都存在开集$U\ni x$与$V\ni y$使$U\cap V=\varnothing$，就称$X$为\term{Hausdorff空间}（Hausdorff space）。
\end{definition}

$\mathbb R$是Hausdorff空间。$\mathbb R^n$也是：若$\bm x\neq\bm y$，它们至少有一个坐标不同，设$x_i\neq y_i$，在第$i$个坐标上取两个不交的开区间，其余坐标取$\mathbb R$，得到两个不交的开矩形。这个论证对一般情形同样适用：两个Hausdorff空间的乘积仍是Hausdorff空间，因为两点在某个因子中不同，就在那个因子中分开。Hausdorff空间的子空间也是Hausdorff空间：若$U,V$在$X$中分开$x,y\in A$，则$U\cap A$与$V\cap A$在$A$中分开它们。所以平面上的圆周、曲线和区域都是Hausdorff空间。

\paragraph{极限的唯一性}
分离条件直接解决了第~\ref{sec:topo-convergence}节留下的极限唯一性问题。

\begin{theorem}[Hausdorff空间中极限唯一]
\label{thm:topo-hausdorff-unique-limit}
在Hausdorff空间中，每个收敛点列只有一个极限。
\end{theorem}
\begin{proof}
设$x_n\to x$且$x_n\to y$，$x\neq y$。取不交的开集$U\ni x$，$V\ni y$。存在$n_1$使$n\geq n_1$时$x_n\in U$，存在$n_2$使$n\geq n_2$时$x_n\in V$。取$n=\max\{n_1,n_2\}$，则$x_n\in U\cap V=\varnothing$，矛盾。
\end{proof}

因此在Hausdorff空间中可以写$\lim_{n\to\infty}x_n=x$。这个定理只需要两个邻域不交；它不要求空间有距离，也不需要可数性。

\paragraph{商空间是否仍可分离}
子空间与乘积保持Hausdorff性质，商空间则不一定，需要逐个检查。先看例~\ref{ex:topo-interval-gluing-circle}的商空间$I/{\sim}$（$0\sim1$）。两个内部点$[s]\neq[t]$（$0<s<t<1$）可以用$(0,1)$中两个不交小区间的像分开，这些像是开集，因为它们的原像就是这两个小区间。内部点$[t]$与粘合点$[0]$之间，取$\varepsilon<\min\{t,1-t\}/2$，用$(t-\varepsilon,t+\varepsilon)$的像与$[0,\varepsilon)\cup(1-\varepsilon,1]$的像分开；后者的原像正是这个并集，所以是开集。于是$I/{\sim}$是Hausdorff空间。

\begin{example}[有两个原点的直线]
\label{ex:topo-line-two-origins}
取两条实轴的并$X=\mathbb R\times\{0,1\}$（$\{0,1\}$取离散拓扑，$X$取乘积拓扑），对每个$t\neq0$令$(t,0)\sim(t,1)$，其余点只与自身等价。商空间$X/{\sim}$看起来像一条实轴，只是原点处有两个点$o_0=[(0,0)]$与$o_1=[(0,1)]$。

这两个原点不能被分开。设$V_0\ni o_0$、$V_1\ni o_1$是开集，则$q^{-1}(V_0)$是含$(0,0)$的开集，包含某个$(-\varepsilon,\varepsilon)\times\{0\}$；同样$q^{-1}(V_1)$包含某个$(-\varepsilon',\varepsilon')\times\{1\}$。取$0<t<\min\{\varepsilon,\varepsilon'\}$，则$[(t,0)]=[(t,1)]$同时属于$V_0$与$V_1$，所以二者相交。由此$X/{\sim}$不是Hausdorff空间，点列$[(1/n,0)]$同时收敛到$o_0$与$o_1$。
\end{example}

这个例子的每一点附近都像实轴的一段，局部看不出任何问题，整体却失去了极限唯一性。所以定义流形时，“局部像欧氏空间”之外还要单独要求Hausdorff条件，见第~\ref{chap:manifold-basics}章。

\subsection{可数性}
\label{sec:topo-countability}

$\mathbb R$有不可数个点，开集更多，但第~\ref{sec:topo-neighborhood-basis}小节已经看到，一个点的全部邻域可以由可数个区间$(x-1/k,x+1/k)$代表。类似地，全部开集能否由可数个基本开集生成？这两个问题分别针对一个点的邻域族和整个空间的开集族。

\begin{definition}[第一可数与第二可数]
\label{def:topo-countability}
若$X$的每个点都有一个可数的邻域基，称$X$\term{第一可数}（first countable）。若$X$有一个可数的拓扑基，称$X$\term{第二可数}（second countable）。
\end{definition}

这里“可数”包括有限。第一可数是逐点的条件，不同点的邻域基可以毫无关系；第二可数是整体的条件，要求一族可数的开集同时服务于所有点。由引理~\ref{lem:topo-basis-to-neighborhood-basis}，从可数拓扑基中取出含$x$的成员，就得到$x$的可数邻域基，所以第二可数蕴含第一可数。

\begin{example}[有理端点区间]
\label{ex:topo-rational-intervals}
在$\mathbb R$中，令$\mathcal B_{\mathbb Q}=\{(p,q):p,q\in\mathbb Q,\ p<q\}$。它是可数集，因为它与$\mathbb Q\times\mathbb Q$的一个子集一一对应，而可数集的有限乘积可数，可参考\BookChapterRef{chap:cardinality}{卷一第3章“基数”}。它是拓扑基：设$U$开，$x\in U$，取$r>0$使$(x-r,x+r)\subseteq U$，由$\mathbb Q$的稠密性取有理数$p\in(x-r,x)$与$q\in(x,x+r)$，则$x\in(p,q)\subseteq U$。

所以$\mathbb R$第二可数。同理，有理端点的开矩形$(p_1,q_1)\times\cdots\times(p_n,q_n)$构成$\mathbb R^n$的可数拓扑基：给定开集$U\ni\bm x$，先取开矩形$\bm x\in R\subseteq U$，再在每个坐标上把区间缩小为含$x_i$的有理端点区间。由于$\mathbb R^n$有不可数个点，这些例子说明，可数性条件约束的是开集族与邻域族的规模，而不是点的个数。
\end{example}

另一方面，第一可数不能推出第二可数。设$X$是不可数集，取离散拓扑。每个点有一个只含一个成员的邻域基$\{\{x\}\}$，所以$X$第一可数。但任何拓扑基$\mathcal B$都必须包含每个单点集：开集$\{x\}$是$\mathcal B$中成员的并，而这些成员都非空且含于$\{x\}$，只能是$\{x\}$本身。于是$\mathcal B$不可数，$X$不第二可数。

第一可数性的意义在于让点列重新可用：若$x\in\overline A$且$x$有可数邻域基$N_1\supseteq N_2\supseteq\cdots$（必要时取前$k$个的交），从每个$N_k\cap A$中选一点，得到的点列就收敛到$x$。例~\ref{ex:topo-cocountable}的余可数拓扑不第一可数，所以点列在那里失效。度量空间都是第一可数的，相应的闭包与闭集序列判据在第~\ref{chap:metric-spaces}章展开。第二可数性则是流形定义中的另一项要求，它保证流形可以由可数个坐标区域拼成。

\subsection{可分性}
\label{sec:topo-separability}

可数性条件谈的是开集族，另一种“可数描述”谈的是点：能否找到可数个点，使它们接近空间中的每一点？在$\mathbb R$中，可数的$\mathbb Q$就做到了这一点，因而数值计算只用有理数也能逼近任何实数。

\begin{definition}[可分空间]
\label{def:topo-separable}
若$X$含有一个可数的稠密子集，称$X$为\term{可分}（separable）空间。
\end{definition}

由例~\ref{ex:topo-rationals-irrationals-dense}，$\mathbb R$是可分的，尽管它不可数；$\mathbb R^n$也可分，$\mathbb Q^n$是可数稠密子集，因为每个非空开集含一个开矩形，开矩形的每个坐标区间都含有理数。可分性与可数性一样，并不限制点的个数。

\paragraph{可分性取决于拓扑}
同一个集合换一个拓扑，可分性可能改变。设$X$是不可数集，取离散拓扑。若$D$稠密，由定理~\ref{thm:topo-dense-open-criterion}，每个非空开集$\{x\}$都与$D$相交，即$x\in D$，所以唯一的稠密子集是$X$本身。$X$不可数，因此不可分。同一个$\mathbb R$，通常拓扑下可分，离散拓扑下不可分；而在平凡拓扑下，任何非空子集都稠密，单点集就够了。

\paragraph{三种可数要求的区别}
到这里出现了三种“可数”：可数稠密点集、可数拓扑基、每点的可数邻域基。它们分别描述点、整体开集族和逐点邻域族，不能互相替代。第二可数能推出可分：在可数拓扑基的每个非空成员中各选一点，得到可数集$D$；每个非空开集都含某个非空基本开集，因而与$D$相交，$D$稠密。

反过来，可分且第一可数的空间也未必第二可数。在$\mathbb R$上以全体半开区间$[a,b)$为基生成拓扑，称为\term{Sorgenfrey直线}（Sorgenfrey line）；两个半开区间的交仍是半开区间或空集，定理~\ref{thm:topo-basis-criterion}的条件成立。每个$[a,b)$都含有理数，所以$\mathbb Q$稠密，空间可分；$\{[x,x+1/k)\}$是$x$的可数邻域基，空间第一可数。但若$\mathcal B$是任意拓扑基，对每个$x$，开集$[x,x+1)$含$x$，必有$B_x\in\mathcal B$使$x\in B_x\subseteq[x,x+1)$，于是$x=\min B_x$。不同的$x$给出不同的$B_x$，所以$\mathcal B$不可数。

在度量空间中，第二可数与可分是等价的，这一结论在第~\ref{chap:metric-spaces}章证明；Sorgenfrey直线说明这个等价依赖距离，并非一般拓扑空间的性质。

\subsection{连通性}
\label{sec:topo-connectedness}

区间$[0,1]$是“一整块”，而$[0,1]\cup[2,3]$由两块组成。描述这种区别有两种思路：一是问空间能否被分成两个互不接触的部分，二是问任意两点能否用一条连续曲线连起来。前者只用开集，称为连通；后者需要显式的路径，称为路径连通。下面分别讨论，并说明在什么局部条件下二者一致。

\subsubsection{连通集}
\label{sec:topo-connected-sets}

$[0,1]\cup[2,3]$可以写成两个不交的非空集合之并，并且每一块都不包含另一块的闭包点，因而两块在这个空间中都是开集。连通就是排除这种分解。

\begin{definition}[连通空间]
\label{def:topo-connected}
若$X$不能写成两个不交的非空开集之并，称$X$\term{连通}（connected）。子集$A\subseteq X$称为连通集，如果$A$作为子空间连通。
\end{definition}

若$X=U\cup V$，$U,V$不交且开，则$U=X\setminus V$也是闭集。所以$X$连通，当且仅当$X$中既开又闭的子集只有$\varnothing$与$X$。这就是第~\ref{sec:topo-closed-sets}小节所说的用途：只要找到一个非平凡的既开又闭子集，空间就被拆开了。离散拓扑下超过一点的空间不连通，平凡拓扑下的空间总是连通的。

\paragraph{实轴上的连通集}
实轴上“没有缺口”的集合是区间：若$x,y\in A$且$x<z<y$，则$z\in A$。下面的定理说明，连通正好刻画了这一点。

\begin{theorem}[实轴上的连通集]
\label{thm:topo-connected-intervals}
子集$A\subseteq\mathbb R$连通，当且仅当$A$是区间。
\end{theorem}
\begin{proof}
若$A$不是区间，存在$x,y\in A$与$z\notin A$，$x<z<y$。则$A=((-\infty,z)\cap A)\cup((z,\infty)\cap A)$，两部分在$A$中相对开、不交，并分别含$x$与$y$，$A$不连通。

设$A$是区间，假设$A=U\cup V$，$U,V$在$A$中相对开、不交且非空。取$a\in U$，$b\in V$，不妨设$a<b$，则$[a,b]\subseteq A$。令$s=\sup(U\cap[a,b])$，则$s\in[a,b]\subseteq A$，下面说明$s$既不能在$U$中，也不能在$V$中。

若$s\in U$，则$s\neq b$，即$s<b$。$U$相对开，存在$\varepsilon>0$使$(s-\varepsilon,s+\varepsilon)\cap A\subseteq U$，于是$(s,\min\{s+\varepsilon,b\})$中的点属于$U\cap[a,b]$，且大于$s$，与上确界的定义矛盾。若$s\in V$，则$s\neq a$，即$s>a$，同样有$\varepsilon>0$使$(s-\varepsilon,s]\cap[a,b]\subseteq V$；这些点都不属于$U$，所以$U\cap[a,b]$中的点都不超过$s-\varepsilon$，与$s$是最小上界矛盾。
\end{proof}

所以$(0,1)$、$[0,1]$、$\mathbb R$都连通，而$(0,1)\cup(2,3)$不连通。证明中用到的是实数的确界原理，这正是“实轴没有缺口”的数学表述；在$\mathbb Q$中同样的论证失败，下文会看到$\mathbb Q$确实不连通。

\paragraph{连续像与介值}
连通性是只用开集表述的性质，连续映射通过原像把开集的分解拉回定义域，因此不能凭空制造分解。

\begin{theorem}[连通集的连续像]
\label{thm:topo-connected-image}
设$f:X\to Y$连续，$X$连通，则$f(X)$连通。
\end{theorem}
\begin{proof}
把$f$看成$X\to f(X)$的映射，由引理~\ref{lem:topo-map-into-subspace}它仍连续。若$f(X)=U\cup V$，$U,V$在$f(X)$中不交、开、非空，则$X=f^{-1}(U)\cup f^{-1}(V)$，两个原像开、不交，且因为$U,V$含有像点而非空，与$X$连通矛盾。
\end{proof}

与定理~\ref{thm:topo-connected-intervals}结合，就得到介值定理的一般形式：若$X$连通，$f:X\to\mathbb R$连续，则$f(X)$是区间；因此若$f$取到值$a<b$，它也取到$a,b$之间的每个值。特别地，同胚保持连通性，所以$[0,1)$与$(0,1)\cup(2,3)$不同胚。

\paragraph{连通分支}
不连通的空间可以分成若干“最大的连通块”。为此需要两个引理，说明连通集在什么情况下合并或扩大后仍连通。

\begin{lemma}[连通集的并与闭包]
\label{lem:topo-connected-union-closure}
\begin{enumerate}
  \item 若一族连通子集$A_i$（$i\in I$）有公共点$p$，则$\bigcup_iA_i$连通；
  \item 若$A$连通，且$A\subseteq B\subseteq\overline A$，则$B$连通。特别地，$\overline A$连通。
\end{enumerate}
\end{lemma}
\begin{proof}
记$S=\bigcup_iA_i$，设$S=U\cup V$，$U,V$在$S$中不交、相对开，不妨$p\in U$。对每个$i$，$A_i\cap U$与$A_i\cap V$在$A_i$中相对开、不交，并覆盖$A_i$；$A_i$连通且$p\in A_i\cap U$，所以$A_i\cap V=\varnothing$，即$A_i\subseteq U$。于是$S\subseteq U$，$V=\varnothing$，第一条成立。

对第二条，设$B=U\cup V$为$B$中不交、相对开的分解。同上，$A$连通，所以$A$整个落在其中一块，不妨$A\subseteq U$。若$V$非空，取$v\in V$，写$V=V'\cap B$，$V'$为$X$的开集。$V'$是$v$的邻域而$v\in\overline A$，所以$V'$与$A$相交；交点属于$A\cap V'\subseteq B\cap V'=V$，与$A\subseteq U$矛盾。故$V=\varnothing$。
\end{proof}

对$x\in X$，令$C(x)$为所有含$x$的连通子集之并。由引理第一条，$C(x)$连通，并且它是含$x$的最大连通子集，称为$x$所在的\term{连通分支}（connected component）。若$C(x)$与$C(y)$相交，二者之并连通且含$x$，由最大性等于$C(x)$，同理等于$C(y)$；所以连通分支两两不交，构成$X$的一个划分。由引理第二条，$\overline{C(x)}$也是含$x$的连通集，故$\overline{C(x)}=C(x)$，连通分支总是闭集。

\begin{example}[连通分支]
\label{ex:topo-components}
$(0,1)\cup(2,3)$的连通分支是$(0,1)$与$(2,3)$。$\mathbb Q$的连通分支是全部单点集：若$A\subseteq\mathbb Q$含两个有理数$p<q$，取无理数$z\in(p,q)$，$(-\infty,z)\cap A$与$(z,\infty)\cap A$就把$A$分开。单点集在$\mathbb Q$中不是开集，所以连通分支一般只是闭集，不必是开集。
\end{example}

\subsubsection{路径连通性}
\label{sec:topo-path-connected}

连通性告诉我们空间不能被拆开，却没有给出从一点到另一点的具体方式。在平面区域中，我们更习惯用一条曲线把两点连起来。下面把“曲线”定义为区间的连续像。

\begin{definition}[路径与路径连通]
\label{def:topo-path-connected}
设$I=[0,1]$。连续映射$\gamma:I\to X$称为从$\gamma(0)$到$\gamma(1)$的\term{路径}（path）。若$X$中任意两点之间都有路径，称$X$\term{路径连通}（path-connected）。
\end{definition}

在$\mathbb R^n$中，线段$\gamma(t)=(1-t)\bm a+t\bm b$是从$\bm a$到$\bm b$的路径，由逐坐标判据它连续；所以$\mathbb R^n$本身、开矩形、开圆盘以及一切凸集都路径连通。

\paragraph{路径分支}
“$x$与$y$之间有路径”是$X$上的等价关系。自反性用常值路径；对称性用反向路径$\bar\gamma(t)=\gamma(1-t)$。传递性需要把从$x$到$y$的路径$\alpha$与从$y$到$z$的路径$\beta$接起来：
\begin{equation}
(\alpha\cdot\beta)(t)=
\begin{cases}
\alpha(2t), & 0\leq t\leq 1/2,\\
\beta(2t-1), & 1/2\leq t\leq 1.
\end{cases}
\label{eq:topo-path-concatenation}
\end{equation}
在$t=1/2$处两段都给出$y$，所以定义无歧义。它的连续性来自下面的粘接引理：若$X=F_1\cup F_2$，$F_1,F_2$闭，$f$在每个$F_j$上的限制连续，则$f$连续。证明只需看闭集$C$的原像：$f^{-1}(C)=(f|_{F_1})^{-1}(C)\cup(f|_{F_2})^{-1}(C)$，每一部分在$F_j$中相对闭，因而形如$F_j$与闭集之交，在$X$中闭。

这个等价关系的等价类称为\term{路径分支}（path component）。路径连通的空间只有一个路径分支。

\paragraph{路径连通蕴含连通}
路径的像是区间$I$的连续像，由定理~\ref{thm:topo-connected-intervals}与定理~\ref{thm:topo-connected-image}是连通集。因此路径连通的空间一定连通：固定$x_0\in X$，对每个$y$取一条从$x_0$到$y$的路径$\gamma_y$，则$X=\bigcup_y\gamma_y(I)$，这些连通集有公共点$x_0$，由引理~\ref{lem:topo-connected-union-closure}，$X$连通。同理，每个路径分支都含于某个连通分支。

\begin{example}[连通但不路径连通]
\label{ex:topo-topologist-sine}
在$\mathbb R^2$中令$G=\{(x,\sin(1/x)):0<x\leq1\}$，$J=\{0\}\times[-1,1]$，$S=G\cup J$。$G$是区间$(0,1]$在连续映射$x\mapsto(x,\sin(1/x))$下的像，因而连通。对$(0,y)\in J$，取$\theta$使$\sin\theta=y$，则$x_k=1/(\theta+2k\pi)\to0$（$k$充分大时$x_k\in(0,1]$），而$(x_k,y)\in G$收敛到$(0,y)$，由定理~\ref{thm:topo-subsequence-closure}，$J\subseteq\overline G$。于是$G\subseteq S\subseteq\overline G$，由引理~\ref{lem:topo-connected-union-closure}，$S$连通。

但没有路径从$(0,0)$连到$(1,\sin1)$。假设$\gamma=(\gamma_1,\gamma_2):I\to S$是这样的路径。$T=\gamma_1^{-1}(\{0\})$是闭集，含$0$而不含$1$，令$t_0=\max T<1$。当$t>t_0$时$\gamma(t)\in G$，所以$\gamma_2(t)=\sin(1/\gamma_1(t))$。由$\gamma_2$在$t_0$处连续，取$\delta\in(0,1-t_0]$使$t\in[t_0,t_0+\delta]$时$\lvert\gamma_2(t)-\gamma_2(t_0)\rvert<1/2$。

记$a=\gamma_1(t_0+\delta)>0$。$\gamma_1$在$[t_0,t_0+\delta]$上从$0$连续变到$a$，由介值定理取遍$(0,a]$。取正整数$k$使$x^+=1/(\pi/2+2k\pi)$与$x^-=1/(3\pi/2+2k\pi)$都小于$a$，则有$t^\pm\in(t_0,t_0+\delta]$使$\gamma_1(t^\pm)=x^\pm$，相应地$\gamma_2(t^+)=1$，$\gamma_2(t^-)=-1$。这两个值相差$2$，却都与$\gamma_2(t_0)$相差小于$1/2$，矛盾。
\end{example}

集合$S$常称为拓扑学家的正弦曲线（topologist's sine curve）。它的问题出在$J$附近：$G$在$J$旁边无限次上下振荡，任何一个小邻域都把$S$切成无穷多段，路径无法在有限时间内穿过。排除这种局部现象，就能让连通与路径连通一致。

\paragraph{局部路径连通}
若对每个$x\in X$及$x$的每个邻域$N$，都存在路径连通的开集$V$使$x\in V\subseteq N$，称$X$\term{局部路径连通}（locally path-connected）。$\mathbb R^n$的开子集都局部路径连通：$\bm x$的邻域含一个以$\bm x$为中心的开矩形，开矩形是凸集，因而路径连通。

\begin{theorem}[局部路径连通时两种分支一致]
\label{thm:topo-locally-path-connected}
设$X$局部路径连通，则$X$的每个路径分支是开集，且路径分支与连通分支相同。特别地，连通且局部路径连通的空间是路径连通的。
\end{theorem}
\begin{proof}
设$P$是一个路径分支，$x\in P$。取路径连通的开集$V\ni x$，$V$中每点经$V$内的路径与$x$相连，所以$V\subseteq P$。于是$P$是开集。$X\setminus P$是其余路径分支之并，也是开集，所以$P$既开又闭。

设$C$是含$P$的连通分支（$P$连通，故含于一个连通分支）。$P\cap C=P$在$C$中既相对开又相对闭，且非空，由$C$连通得$P=C$。最后一句是$X$只有一个连通分支的情形。
\end{proof}

于是$\mathbb R^n$中连通的开集都是路径连通的，连通性可以用画出一条曲线来核验。例~\ref{ex:topo-topologist-sine}中的$S$在$J$上的点处不局部路径连通，定理的条件不满足。流形在每点附近同胚于欧氏空间的开集，因而局部路径连通，第~\ref{chap:manifold-basics}章将用这一点把流形的连通与路径连通等同起来。

\subsection{紧性}
\label{sec:topo-compactness}

闭区间上的连续函数有最大值，开区间上则未必；卷一把这一区别归结为闭区间的“有限子覆盖”性质。这个性质只用开集表述，因此可以搬到任何拓扑空间：它要求无论用什么方式以局部区域覆盖整个对象，都只需其中有限个。本小节先给出定义，再讨论闭包紧的集合，最后说明连续映射怎样保持紧性。

\subsubsection{紧性的定义}
\label{sec:topo-compact-definition}

一族开集$\{U_i\}_{i\in I}$若满足$\bigcup_iU_i=X$，称为$X$的\term{开覆盖}（open cover）；从中选出的、仍覆盖$X$的子族称为子覆盖，有限个成员组成的子覆盖称为\term{有限子覆盖}（finite subcover）。

\begin{definition}[紧空间与紧集]
\label{def:topo-compact}
若$X$的每个开覆盖都有有限子覆盖，称$X$\term{紧致}（compact），也称紧空间。子集$K\subseteq X$称为紧集，如果$K$作为子空间是紧的。
\end{definition}

子空间的开集是$U\cap K$的形式，所以$K$是紧集，当且仅当每一族满足$K\subseteq\bigcup_iU_i$的$X$的开集都有有限个成员仍覆盖$K$。以后检验紧集时直接用环境空间的开集，不必先转成相对开集。

\paragraph{某一个覆盖与所有覆盖}
定义中的量词是“每个开覆盖”。$(0,1)$有一个只含一个成员的开覆盖$\{(0,1)\}$，这不说明任何问题。考虑覆盖$\{(1/n,1):n\geq2\}$：每点$x\in(0,1)$在$n>1/x$时属于$(1/n,1)$，所以它覆盖$(0,1)$；但有限个成员之并是$(1/N,1)$，漏掉了$(0,1/N]$。所以$(0,1)$不紧。同理，$\{(-n,n):n\geq1\}$说明$\mathbb R$不紧。

闭区间$[a,b]$则是紧集，卷一已经用共同半径的论证证明了这一点。由下面的引理，$\mathbb R$中有界闭集都是紧集：它是某个$[a,b]$的闭子集。

\begin{lemma}[紧空间的闭子集]
\label{lem:topo-closed-in-compact}
紧空间$X$的闭子集$F$是紧集。
\end{lemma}
\begin{proof}
设$\{U_i\}$是$X$的开集，覆盖$F$。加上开集$X\setminus F$，得到$X$的开覆盖，取出有限子覆盖。去掉其中的$X\setminus F$（若出现），剩下的有限个$U_i$仍覆盖$F$。
\end{proof}

其他例子：有限集总是紧的，因为每点只需选一个含它的成员。无限集取离散拓扑时不紧，单点集组成的开覆盖没有有限子覆盖。平凡拓扑下的空间总是紧的，因为唯一的非空开集就是全空间。最后一例也说明，紧性本身不要求点能被分开。

在实轴上，紧性还可以用“每个点列有收敛子列”来刻画，卷一的证明正是在两种说法之间往返。一般拓扑空间中，两种说法可以不等价，点列也未必能探测拓扑，所以紧性以开覆盖定义。在度量空间中二者等价，并且紧等价于完备且全有界，这些判据在第~\ref{chap:metric-spaces}章建立。

\subsubsection{相对紧集}
\label{sec:topo-relatively-compact}

$(0,1)$本身不紧，但它离紧只差端点：补上闭包点后得到$[0,1]$，是紧集。很多估计只需要集合落在某个紧集里，而不需要集合本身紧，这引出下面的概念。

\begin{definition}[相对紧集]
\label{def:topo-relatively-compact}
设$A\subseteq X$。若$A$在$X$中的闭包$\overline A$是紧集，称$A$在$X$中\term{相对紧}（relatively compact）。
\end{definition}

$(0,1)$在$\mathbb R$中相对紧，但它既不闭也不紧；所以相对紧是比紧弱的条件。$\mathbb R$的每个有界子集都相对紧，因为它的闭包是有界闭集。紧集是否一定相对紧，涉及紧集的闭包；在Hausdorff空间中，这个问题由下面的引理回答。

\begin{lemma}[Hausdorff空间中的紧集是闭集]
\label{lem:topo-compact-closed-hausdorff}
设$X$是Hausdorff空间，$K\subseteq X$是紧集，则$K$是闭集。
\end{lemma}
\begin{proof}
取$x\notin K$。对每个$y\in K$，由Hausdorff条件取不交的开集$U_y\ni x$与$V_y\ni y$。$\{V_y\}_{y\in K}$覆盖$K$，取有限个$V_{y_1},\dots,V_{y_m}$覆盖$K$。则$U=U_{y_1}\cap\cdots\cap U_{y_m}$是含$x$的开集，与$V_{y_1}\cup\cdots\cup V_{y_m}\supseteq K$不交。所以$X\setminus K$是它每个点的邻域，是开集。
\end{proof}

由此，在Hausdorff空间中，紧集等于自己的闭包，因而相对紧；任何含于紧集$K$的子集$A$也相对紧，因为$\overline A$是$K$的闭子集，由引理~\ref{lem:topo-closed-in-compact}紧。没有Hausdorff条件时引理可以失效：平凡拓扑下超过一点的空间里，每个单点集都紧，却不闭。

\paragraph{相对紧依赖环境空间}
同一个集合在不同的环境中可能一个相对紧、一个不相对紧。把$(0,1)$看作空间$X=(0,1)$自身的子集，它在$X$中的闭包就是$X$，而$X$不紧，所以$(0,1)$在$(0,1)$中不相对紧；在$\mathbb R$中，它的闭包多了两个端点，成了紧集。这说明相对紧描述的是集合与环境的关系：环境中是否有足够的点让集合“收拢”。在度量空间中，相对紧可以用点列判据与全有界性来检验，这留到第~\ref{chap:metric-spaces}章。

\subsubsection{紧集的连续像}
\label{sec:topo-compact-images}

闭区间上连续函数取到最大值，这一结论的拓扑核心是：连续映射把紧集映成紧集，而实轴上的紧集有最大元。

\begin{theorem}[紧集的连续像]
\label{thm:topo-compact-image}
设$f:X\to Y$连续，$K\subseteq X$是紧集，则$f(K)$是$Y$中的紧集。
\end{theorem}
\begin{proof}
设$\{V_i\}$是$Y$的开集，覆盖$f(K)$。则$\{f^{-1}(V_i)\}$是$X$的开集，覆盖$K$：每个$x\in K$的像$f(x)$属于某个$V_i$。由$K$紧，取有限个$f^{-1}(V_{i_1}),\dots,f^{-1}(V_{i_m})$覆盖$K$，则$V_{i_1},\dots,V_{i_m}$覆盖$f(K)$。
\end{proof}

证明的方向值得注意：开覆盖在输出空间给出，经连续性拉回到输入空间，在那里用紧性有限化，再推回输出空间。连续性负责“拉回后仍是开覆盖”，紧性负责“有限化”，二者缺一不可。

\begin{theorem}[最大值与最小值]
\label{thm:topo-extreme-values}
设$K$是非空紧空间，$f:K\to\mathbb R$连续，则$f$有界，并且取到最大值与最小值。
\end{theorem}
\begin{proof}
由定理~\ref{thm:topo-compact-image}，$f(K)$是$\mathbb R$中的非空紧集。开区间$(-n,n)$（$n\geq1$）覆盖$f(K)$，有限个就够，所以$f(K)$含于某个$(-n,n)$，$f$有界。$\mathbb R$是Hausdorff空间，由引理~\ref{lem:topo-compact-closed-hausdorff}，$f(K)$是闭集。令$M=\sup f(K)$，对每个$\varepsilon>0$，$(M-\varepsilon,M+\varepsilon)$含$f(K)$中的点，所以$M\in\overline{f(K)}=f(K)$，即最大值$M$被取到。对$-f$用同样的论证得到最小值。
\end{proof}

两个条件都不能去掉。$f(x)=x$在$(0,1)$上连续，但$(0,1)$不紧，$f$没有最大值。在紧集$[0,1]$上，令$f(x)=x$（$x<1$）且$f(1)=0$，则$f$不连续，上确界$1$取不到。

\paragraph{同胚保持紧性}
若$h:X\to Y$是同胚，$X$紧，则$Y=h(X)$紧；反过来用$h^{-1}$。所以紧性与连通性一样是拓扑性质。由此，$[0,1]$与$(0,1)$不同胚，尽管二者都连通。紧性还给出判断同胚的一个有力工具，它回答了第~\ref{sec:topo-quotients}小节留下的问题。

\begin{theorem}[从紧空间到Hausdorff空间的连续双射]
\label{thm:topo-compact-hausdorff-homeomorphism}
设$X$紧，$Y$是Hausdorff空间，$f:X\to Y$是连续双射，则$f$是同胚。
\end{theorem}
\begin{proof}
只需证明$f^{-1}$连续，由定理~\ref{thm:topo-continuity-equivalent}，即证明$f$把$X$的闭集映成闭集，因为$(f^{-1})^{-1}(F)=f(F)$。设$F$在$X$中闭，由引理~\ref{lem:topo-closed-in-compact}，$F$紧；由定理~\ref{thm:topo-compact-image}，$f(F)$紧；由引理~\ref{lem:topo-compact-closed-hausdorff}，$f(F)$在$Y$中闭。
\end{proof}

现在回到例~\ref{ex:topo-interval-gluing-circle}。商空间$I/{\sim}=q(I)$是紧集$I$的连续像，因而紧；$S^1$是Hausdorff空间$\mathbb R^2$的子空间，因而是Hausdorff空间。所以连续双射$\bar p:I/{\sim}\to S^1$是同胚，“把线段两端粘起来”确实得到圆周。对照例~\ref{ex:topo-interval-to-circle}，$[0,1)\to S^1$同样是连续双射，却不是同胚；定理在那里不适用，因为$[0,1)$不紧。
